\documentclass[10pt,a4paper]{amsart}
\usepackage[margin=19mm]{geometry}
\usepackage{amsmath,amssymb,mathtools}
\usepackage[scr=rsfs,cal=euler]{mathalfa}
\usepackage{xfrac}
\usepackage[unicode,hidelinks,bookmarks=false]{hyperref}
\newcommand{\ie}{i.e.\ }
\newcommand{\cf}{cf.\ }
\newcommand{\resp}{respectively\ }
\newcommand{\Subsection}{\subsection}
\providecommand{\nobreakdash}{\nobreak\hbox{-}}
\setlength{\emergencystretch}{2em}
\multlinegap0pt
\usepackage{bm}
\usepackage{bbm}
\usepackage{dsfont}
\usepackage{xspace}
\usepackage{cancel}
\usepackage{mathtools}
\usepackage{amsthm}
\makeatletter
\@ifundefined{theorem}{\newtheorem{theorem}{Theorem}}{}
\makeatother
\newcommand{\E}{\mathrm{E}}
\newcommand{\R}{{\mathbb{R}}}
\newcommand{\e}{\mathrm{e}}
\renewcommand{\d}{\mathrm{d}}
\title[Spatial fluctuation for stochastic heat equation with H\"old]{Spatial fluctuation for stochastic heat equation with H\"older coefficients}
\author{Carl Mueller, Fei Pu}
\date{Source: arXiv:2510.00807v1 (CC BY 4.0)}
\begin{document}
\maketitle

\noindent\emph{Source and licence.} Spatial fluctuation for stochastic heat equation with H\"older coefficients. Version arXiv:2510.00807v1 (CC BY 4.0), distributed under the Creative Commons Attribution 4.0 International licence, as recorded in the arXiv licence metadata for this exact version and shown on the abstract page https://arxiv.org/abs/2510.00807v1. Full rights notice: licenses/arxiv-2510.00807-CC-BY-4.0.txt.\par\smallskip

\noindent\textit{Editorial scope.} Counted unit: the Theorem whose source label is \texttt{th:LDP} in the article (source0.tex lines 496--501, source label \texttt{th:LDP}), delivered together with the article's own complete original proof of it (source0.tex lines 502--629, one \texttt{proof} environment of 128 lines). No mathematics is written, repaired or re-derived: statement and proof are the article's own text line for line. Required context retained uncounted: eq:SHE (lines 89--94). Every definition, display and label of those lines is retained unchanged. Omitted from this package and not redistributed: 1--88, 95--495, 630--1098. Nothing inside a retained excerpt is omitted or altered: every retained body is delivered exactly as the article's lines are written, so no omission receipt of any kind accompanies this package. Citations: the retained excerpts carry no citation command at all, so no bibliography entry of the article is transcribed and no reference is redistributed. Reference adaptation: none was needed and none was made; every reference inside the delivered unit resolves inside it, so no unresolved reference or citation is produced. Printed slips kept literal: no printed word, symbol, hypothesis or number of the counted statement or of its proof was altered. Layout and class: the article's own class is \texttt{amsart} with packages including graphicx, xcolor, cite, mathtools, amssymb, amsmath, amsthm, mathrsfs, paralist, bm, esint, setspace, hyperref, color; the wrapper is the lane's pinned \texttt{amsart} editorial wrapper at 10pt on A4, which restates the theorem-like environments the retained text needs and adds the article's own macro definitions that the retained text needs (\texttt{\textbackslash E}, \texttt{\textbackslash R}, \texttt{\textbackslash d}, \texttt{\textbackslash e}), with extra wrapper packages (bm, bbm, dsfont, xspace, cancel, mathtools). The delivered numbering is the unit's own. Assets: no retained span contains a figure environment, a graphics-inclusion command, a PDF-page inclusion, a table or a TikZ picture; no image file is copied into this package, the package declares no asset dependency and no image file is redistributed with it. Licence: version arXiv:2510.00807v1 of this article is distributed under the Creative Commons Attribution 4.0 International licence, as recorded in the arXiv licence metadata for this exact version and shown on the abstract page https://arxiv.org/abs/2510.00807v1. A full rights notice is delivered at licenses/arxiv-2510.00807-CC-BY-4.0.txt. Characters: every retained excerpt is pure ASCII, so the delivered engine, XeLaTeX with its default Latin Modern text font, reports no missing character.
\begin{align}\label{eq:SHE}
\begin{cases}
\partial_tu(t,x)=\frac12\partial_x^2u(t,x)+ u(t,x)^\gamma \xi(t,x), \quad t>0, x\in \R\\
u(0)\equiv1,
\end{cases}
\end{align}

\begin{theorem}\label{th:LDP}
          Assume $\gamma=\frac12$. Fix $t>0$. For $\lambda>0$, it holds that
          \begin{align}\label{eq:LDP2}
           \lim_{N\to\infty} \frac1N\log \E\left[\e^{-\lambda S_{N,t}}\right]= -\frac{2\lambda}{\lambda t+2}.
          \end{align}
\end{theorem}

\begin{proof}
           When $\gamma=\frac12$, the solution to \eqref{eq:SHE} corresponds to the density of super Brownian motion. By 
           the exponential duality, we write for $\lambda>0$
\begin{align*}
\E\left[\e^{-\lambda S_{N,t}}\right] =\E\left[\e^{-\lambda \langle u(t, \cdot)\,, 
\bm{1}_{[0, N]}\rangle}\right] = \e^{-\langle 1, v(t, \cdot)\rangle},
\end{align*}
where $v$ solves the heat equation
\begin{align}\label{eq:HE1}
\begin{cases}
\partial_tv= \frac12 \partial_x^2 v -\frac12 v^2,\\
v(0, x)= \lambda \bm{1}_{[0, N]}(x).
\end{cases}
\end{align}
Thus, it remains to verify 
\begin{align}\label{eq:LDPlim}
\lim_{N\to\infty} \frac{\int_\R v(t,x)\, \d x}{N} = \frac{2\lambda}{\lambda t+2}.
\end{align}


The solution to \eqref{eq:HE1} is nonnegative and satisfies the following integral equation
\begin{align*}
v(t,x)= \int_\R \lambda \bm{1}_{[0, N]}(y) p_t(x-y)\d y -\frac12 \int_0^t\int_\R p_{t-s}(x-y)v(s, y)^2\d y\d s
\end{align*}
and hence we have 
\begin{align}\label{eq:vbound}
v(t,x ) \leq  \int_\R \lambda \bm{1}_{[0, N]}(y) p_t(x-y)\d y, \quad t>0, x\in\R. 
\end{align}

Let $w$ be the solution to 
\begin{align}\label{eq:HE2}
\begin{cases}
\partial_tw= \frac12 \partial_x^2 w -\frac12 w^2,\\
w(0, x)= \lambda.
\end{cases}
\end{align}
Note that $w$ depends only on $t$, so will write $w(t)$. Thus, $w$ solves an ODE with an explicit formula given by
\begin{align*}
w(t)= \frac{2\lambda}{\lambda t+2}, \quad t\geq0.
\end{align*}
Let 
\begin{align*}
h(t,x)= w(t,x)-v(t,x).
\end{align*}
Because the initial value of $w$ is larger than that of $v$, by the comparison principle, $h$ is nonnegative (and bounded above by $\lambda$). In addition,  we see that $h$ satisfies
\begin{align*}
\partial_t h &=\partial_t w-\partial_t v
=\frac12\partial_x^2 w-\frac12w^2-\frac12\partial_x^2v+\frac12v^2\\
&=\frac12\partial_x^2h-\frac12(w+v)h \leq \frac12\partial_x^2h.
\end{align*}
%By the maximum principle for the Cauchy problem (see \cite[p.57, Theorem 6]{Eva10}), we obtain
A supersolution for $h$ is $\bar{h}$, where
\begin{align}\label{eq:h}
h\leq \bar{h},
\end{align}
and $\bar{h}$ solves
\begin{align}\label{eq:HE3}
\begin{cases}
\partial_t \bar{h}=\frac12\partial_x^2\bar{h},\\
\bar{h}(0, x)= \lambda \bm{1}_{[0, N]^c}(x).
\end{cases}
\end{align}
The solution to \eqref{eq:HE3} is given by
\begin{align*}
\bar{h}(t,x)= \int_\R \lambda \bm{1}_{[0, N]^c}(y)p_t(x-y)\d y.
\end{align*}


In order to prove \eqref{eq:LDPlim},
we use \eqref{eq:vbound} and \eqref{eq:h} to see that
\begin{align*}
&\left|\frac1N\int_\R v(t,x)\d x -w(t)\right| 
\leq \frac1N\left[\int_{[0, N]^c} v(t,x)\d x  + \int_0^N \bar{h}(t,x)\d x\right]\\
&\qquad\leq  \frac\lambda N\Big[\int_{[0, N]^c}\left(  \int_\R \bm{1}_{[0, N]}(y) p_t(x-y)\d y\right)\d x  \\
&\qquad\qquad\qquad\qquad \qquad\qquad
+ \int_0^N \left( \int_\R \bm{1}_{[0, N]^c}(y)p_t(x-y)\d y\right)\d x\Big]\\
&\qquad =\frac{2\lambda} N\int_{[0, N]^c}\left(  \int_\R \bm{1}_{[0, N]}(y) p_t(x-y)\d y\right)\d x.
\end{align*}
We only need to show that
\begin{align}\label{eq:lim2}
\lim_{N\to\infty} \frac1N\int_{-\infty}^0 \left(\int_0^N p_t(x-y)\d y\right) \d x =0
\end{align}
and
\begin{align}\label{eq:lim}
  \lim_{N\to\infty} \frac1N\int^{\infty}_N \left(\int_0^N p_t(x-y)\d y\right) \d x =0
\end{align}
For the integral in \eqref{eq:lim2}, we write
\begin{align*}
 \frac1N\int_{-\infty}^0 \left(\int_0^N p_t(x-y)\d y\right) \d x= \frac1N \int_0^N \left(\int_{-\infty}^{-y}p_t(z)\d z\right)\d y.
\end{align*}
By L'H\^opital's rule, 
\begin{align*}
\lim_{N\to\infty} \frac1N \int_0^N \left(\int_{-\infty}^{-y}p_t(z)\d z\right)\d y = \lim_{N\to\infty} \int_{-\infty}^{-N}p_t(z)\d z =0.
\end{align*}
For the  integral in \eqref{eq:lim}, we write
\begin{align*}
&\frac1N\int^{\infty}_N \left(\int_0^N p_t(x-y)\d y\right) \d x\\
& \quad= \frac1N\int^{\infty}_{2N} \left(\int_0^N p_t(x-y)\d y\right) \d x
+ \frac1N\int^{2N}_N \left(\int_0^N p_t(x-y)\d y\right) \d x.
\end{align*}
It is clear that
\begin{equation*} %\label{eq:clear}
\begin{split}
 \frac1N\int^{\infty}_{2N} \left(\int_0^N p_t(x-y)\d y\right) \d x &\leq \frac1N \int^{\infty}_{2N} \left(\int_0^N p_t(x/2)\d y\right) \d x\\
 &=\int_{2N}^{\infty}p_t(x/2)\d x \to 0, \quad \text{as $N\to\infty$.}
\end{split}
\end{equation*}
To justify the first line, we briefly explain why $|x-y|\ge x/2$ for $x,y$ in 
the domain of integration.  
Indeed, for such $x,y$ we have $x\ge2N$ and $y\le N$.  Therefore $y\le x/2$ 
and so $|x-y|=x-y\ge x-(x/2)=x/2$.  
Moreover, by a change of variable, 
\begin{align*}
\frac1N\int^{2N}_N \left(\int_0^N p_t(x-y)\d y\right) \d x&=\int_1^2\d x \int_0^1 \d y \frac{N}{\sqrt{2\pi t}} \e^{-\frac{N^2(x-y)^2}{2t}}\\
&=\int_1^2\d x \int_0^1 \d y \, \frac1{x-y}\frac{N(x-y)}{\sqrt{2\pi t}} \e^{-\frac{N^2(x-y)^2}{2t}}.
\end{align*}
Since 
\begin{align*}
\sup_{N\geq1, x, y\in \R} \left[N|x-y|\e^{-\frac{N^2|x-y|^2}{2t}}\right]<\infty \quad \text{and}\quad \int_1^2\d x \int_0^1 \d y \, \frac1{x-y}<\infty,
\end{align*}
by the dominated convergence theorem, 
\begin{align*}
\lim_{N\to\infty}\frac1N\int^{2N}_N \left(\int_0^N p_t(x-y)\d y\right) \d x=0.
\end{align*}
Therefore,  \eqref{eq:lim} holds and it proves \eqref{eq:LDPlim}. 

The proof of Theorem \ref{th:LDP} is complete.
\end{proof}

\end{document}
